% GENERATED FILE. Edit canonical lessons, then run npm run generate:books.
% canonical-source-hash: fnv1a64:faae1d9997c6bdd6
% canonical-lessons: GU-C393-enjin, GU-C393-lain, GU-C393-janksan, GU-C393-siti, GU-C393-jodvu

\chapter{Engines and Lines}
\label{ch:gu-engines-and-lines}
\hlchaptermodality{393}

\begin{chapteropening}
\textbf{By the end of this chapter:} \emph{I can talk about the engine, the line and the junction.}

Complete the last lesson of chapter 393: I can talk about the engine, the line and the junction.
\end{chapteropening}

\section[\texorpdfstring{\gu{એન્જિન} (enjin)}{એન્જિન (enjin)}]{\gu{એન્જિન} (enjin) — an engine (of a train)}

\label{lesson:GU-C393-enjin}



\begin{quote}
Before the new one: say the Gujarati for to pester, then the Gujarati for to hope.
\end{quote}

\subsection*{You'll want to know: \gu{એન્જિન}}
\textbf{\gu{એન્જિન}} — \emph{enjin} — \textquotedblleft{}an engine (of a train)\textquotedblright{}.

\begin{quote}
\gu{આ} \gu{ટ્રેનનું} \gu{એન્જિન} \gu{ઘણું} \gu{જૂનું} \gu{છે}.
\end{quote}

\begin{quote}
\emph{ā ṭrennũ enjin ghaṇũ jūnũ che.}
\end{quote}

\begin{quote}
The engine of this train is very old.
\end{quote}

A neuter noun.

\begin{cousinweb}
\textbf{\gu{એન્જિન}} (\emph{enjin}) is borrowed from English \emph{engine}; on the railway it is the engine at the front that pulls the train.
\end{cousinweb}

\begin{grammarlens}[title={What you've built}]
One more word for getting around.
\end{grammarlens}

\subsection*{Your turn}
\begin{itemize}
\raggedright
  \item \emph{Say it:} \emph{enjin}
  \item \emph{Say it:} \emph{enjin}, once more
  \item \emph{Say it:} \emph{ā ṭrennũ enjin ghaṇũ jūnũ che.} — The engine of this train is very old.
  \item \emph{Recall:} say the Gujarati for to pester, then the Gujarati for to hope, then say \emph{enjin} again
\end{itemize}

\begin{tcolorbox}[breakable,colback=teal!4,colframe=teal!35!black,title={Before you move on}]
What does \gu{એન્જિન} mean? (\textquotedblleft{}An engine (of a train)\textquotedblright{}.) Say it once more.
\end{tcolorbox}
\section[\texorpdfstring{\gu{લાઈન} (lāīn)}{લાઈન (lāīn)}]{\gu{લાઈન} (lāīn) — a railway line; a line}

\label{lesson:GU-C393-lain}



\begin{quote}
Before the new one: say the Gujarati for to hope, then the Gujarati for an engine (of a train).
\end{quote}

\subsection*{You'll want to know: \gu{લાઈન}}
\textbf{\gu{લાઈન}} — \emph{lāīn} — \textquotedblleft{}a railway line; a line\textquotedblright{}.

\begin{quote}
\gu{અમારા} \gu{ગામ} \gu{પાસે} \gu{રેલવેની} \gu{લાઈન} \gu{છે}.
\end{quote}

\begin{quote}
\emph{amārā gām pāse relvenī lāīn che.}
\end{quote}

\begin{quote}
There is a railway line near our village.
\end{quote}

A feminine noun.

\begin{cousinweb}
\textbf{\gu{લાઈન}} (\emph{lāīn}) is borrowed from English \emph{line}, here the line of track a train runs on.
\end{cousinweb}

\begin{grammarlens}[title={What you've built}]
One more word for getting around.
\end{grammarlens}

\subsection*{Your turn}
\begin{itemize}
\raggedright
  \item \emph{Say it:} \emph{lāīn}
  \item \emph{Say it:} \emph{lāīn}, once more
  \item \emph{Say it:} \emph{amārā gām pāse relvenī lāīn che.} — There is a railway line near our village.
  \item \emph{Recall:} say the Gujarati for to hope, then the Gujarati for an engine (of a train), then say \emph{lāīn} again
\end{itemize}

\begin{tcolorbox}[breakable,colback=teal!4,colframe=teal!35!black,title={Before you move on}]
What does \gu{લાઈન} mean? (\textquotedblleft{}A railway line; a line\textquotedblright{}.) Say it once more.
\end{tcolorbox}
\section[\texorpdfstring{\gu{જંકશન} (jankśan)}{જંકશન (jankśan)}]{\gu{જંકશન} (jankśan) — a junction (of railway lines)}

\label{lesson:GU-C393-janksan}



\begin{quote}
Before the new one: say the Gujarati for an engine (of a train), then the Gujarati for a railway line.
\end{quote}

\subsection*{You'll want to know: \gu{જંકશન}}
\textbf{\gu{જંકશન}} — \emph{jankśan} — \textquotedblleft{}a junction (of railway lines)\textquotedblright{}.

\begin{quote}
\gu{આ} \gu{જંકશન} \gu{પર} \gu{ત્રણ} \gu{લાઈન} \gu{મળે} \gu{છે}.
\end{quote}

\begin{quote}
\emph{ā jankśan par traṇ lāīn maḷe che.}
\end{quote}

\begin{quote}
Three lines meet at this junction.
\end{quote}

A neuter noun.

\begin{cousinweb}
\textbf{\gu{જંકશન}} (\emph{jankśan}) is borrowed from English \emph{junction}, a station where railway lines meet.
\end{cousinweb}

\begin{grammarlens}[title={What you've built}]
One more word for getting around.
\end{grammarlens}

\subsection*{Your turn}
\begin{itemize}
\raggedright
  \item \emph{Say it:} \emph{jankśan}
  \item \emph{Say it:} \emph{jankśan}, once more
  \item \emph{Say it:} \emph{ā jankśan par traṇ lāīn maḷe che.} — Three lines meet at this junction.
  \item \emph{Recall:} say the Gujarati for an engine (of a train), then the Gujarati for a railway line, then say \emph{jankśan} again
\end{itemize}

\begin{tcolorbox}[breakable,colback=teal!4,colframe=teal!35!black,title={Before you move on}]
What does \gu{જંકશન} mean? (\textquotedblleft{}A junction (of railway lines)\textquotedblright{}.) Say it once more.
\end{tcolorbox}
\section[\texorpdfstring{\gu{સીટી} (sīṭī)}{સીટી (sīṭī)}]{\gu{સીટી} (sīṭī) — a whistle}

\label{lesson:GU-C393-siti}



\begin{quote}
Before the new one: say the Gujarati for a railway line, then the Gujarati for a junction (of railway lines).
\end{quote}

\subsection*{You'll want to know: \gu{સીટી}}
\textbf{\gu{સીટી}} — \emph{sīṭī} — \textquotedblleft{}a whistle\textquotedblright{}.

\begin{quote}
\gu{ટ્રેનની} \gu{સીટી} \gu{સાંભળીને} \gu{બાળકો} \gu{દોડ્યાં}.
\end{quote}

\begin{quote}
\emph{ṭrennī sīṭī sā̃bhḷīne bāḷko doḍyā̃.}
\end{quote}

\begin{quote}
Hearing the train's whistle, the children ran.
\end{quote}

A feminine noun.

\begin{cousinweb}
\textbf{\gu{સીટી}} (\emph{sīṭī}) goes back to \emph{*sīṭṭa}, \textquotedblleft{}a whistle\textquotedblright{}, a form CDIAL reconstructs and treats as imitative of the sound itself; CDIAL lists the Gujarati word (CDIAL 13427).
\end{cousinweb}

\begin{grammarlens}[title={What you've built}]
One more word for getting around.
\end{grammarlens}

\subsection*{Your turn}
\begin{itemize}
\raggedright
  \item \emph{Say it:} \emph{sīṭī}
  \item \emph{Say it:} \emph{sīṭī}, once more
  \item \emph{Say it:} \emph{ṭrennī sīṭī sā̃bhḷīne bāḷko doḍyā̃.} — Hearing the train's whistle, the children ran.
  \item \emph{Recall:} say the Gujarati for a railway line, then the Gujarati for a junction (of railway lines), then say \emph{sīṭī} again
\end{itemize}

\begin{tcolorbox}[breakable,colback=teal!4,colframe=teal!35!black,title={Before you move on}]
What does \gu{સીટી} mean? (\textquotedblleft{}A whistle\textquotedblright{}.) Say it once more.
\end{tcolorbox}
\section[\texorpdfstring{\gu{જોડવું} (joḍvũ)}{જોડવું (joḍvũ)}]{\gu{જોડવું} (joḍvũ) — to couple (a carriage or an engine); to attach}

\label{lesson:GU-C393-jodvu}



\begin{quote}
Before the new one: say the Gujarati for a junction (of railway lines), then the Gujarati for a whistle.
\end{quote}

\subsection*{You'll want to know: \gu{જોડવું}}
\textbf{\gu{જોડવું}} — \emph{joḍvũ} — \textquotedblleft{}to couple (a carriage or an engine); to attach\textquotedblright{}.

\begin{quote}
\gu{સ્ટેશન} \gu{પર} \gu{તેમણે} \gu{ટ્રેનમાં} \gu{નવું} \gu{એન્જિન} \gu{જોડ્યું}.
\end{quote}

\begin{quote}
\emph{sṭeśan par temṇe ṭrenmā̃ navũ enjin joḍyũ.}
\end{quote}

\begin{quote}
At the station they coupled a new engine to the train.
\end{quote}

Past, with \gu{તેણે} (he or she …): \gu{તેણે} \gu{એન્જિન} \gu{જોડ્યું}, \emph{teṇe enjin joḍyũ}.

\begin{cousinweb}
\textbf{\gu{જોડવું}} (\emph{joḍvũ}) goes back to \emph{*yuṭati}, \textquotedblleft{}is joined\textquotedblright{}, a form CDIAL reconstructs; CDIAL lists the Gujarati verb, \textquotedblleft{}to join, to fasten\textquotedblright{} (CDIAL 10496). \gu{જોડી}, a pair, and \gu{જોડાવું}, to join, are its relatives.
\end{cousinweb}

\begin{grammarlens}[title={What you've built}]
One more word for getting around.
\end{grammarlens}

\subsection*{Your turn}
\begin{itemize}
\raggedright
  \item \emph{Say it:} \emph{joḍvũ}
  \item \emph{Say it:} \emph{joḍvũ}, once more
  \item \emph{Say it:} \emph{sṭeśan par temṇe ṭrenmā̃ navũ enjin joḍyũ.} — At the station they coupled a new engine to the train.
  \item \emph{Recall:} say the Gujarati for a junction (of railway lines), then the Gujarati for a whistle, then say \emph{joḍvũ} again
\end{itemize}

\begin{tcolorbox}[breakable,colback=teal!4,colframe=teal!35!black,title={Before you move on}]
What does \gu{જોડવું} mean? (\textquotedblleft{}To couple (a carriage or an engine); to attach\textquotedblright{}.) Say it once more.
\end{tcolorbox}
